\section{Results: function-aligned RL moves the diversity/fidelity frontier}
\label{sec:res:frontier}

\plan{\textbf{Claim.} On an established public benchmark, \cleo{} produces
libraries with the sequence coverage of high-temperature LigandMPNN and the
pass rate of low-temperature LigandMPNN --- a combination no single temperature
achieves.

\textbf{Why this section leads.} The PETase result rests on assays that take
months and cannot be re-run by a referee. This one is fully computational, uses
a published benchmark and a published deposit, and can be reproduced end to end.
It buys the method credibility before any wet-lab claim is weighed.

\textbf{Status: \statusscaffold{}}, and close to \statusdone{}. Both the \cleo{}
arms and the full temperature ladder are measured. What remains: reconcile the
two \cleo{} rows below (Block 3 quotes the corrected 200-step anchored arm; the
notebook's E10 table still carries a superseded pre-fix peak window), rebuild
the two PCA panels, fold in the E20 rescue result, and attach the run-to-run
error bars from E21.}

\subsection*{Block 1 --- Setup and the unit of comparison}

We evaluate on the atomized motif enforcement (AME) benchmark introduced with
RFdiffusion2, comprising 41 catalytic sites scaffolded into 100 backbones
each. The published protocol assigns 8 LigandMPNN sequences per backbone and
predicts each with 5 Chai models; a prediction succeeds when the motif
all-atom RMSD to the design is below \SI{1.5}{\angstrom} and no backbone atom
clashes with the ligand.

\plan{The aggregation rule matters more than it looks and is easy to get wrong.
We verified it against the deposit rather than assuming: the per-site number
reported in \texttt{AME\_benchmark\_summary.csv} is the fraction of the 100
backbones with \emph{at least one} passing prediction out of the 40, matching
to zero error across all 41 sites and both methods. It is a best-of-40, not a
mean over predictions. Scoring one prediction per sequence and comparing rates
therefore understates a comparable method by roughly threefold --- in the
deposit, only 36\,\% of the sequences that pass on some Chai model pass on all
five. Every number below is on their unit.}

\begin{planlist}
  \item Pilot backbones: three drawn from the \texttt{pass} class, at
        32/40, 28/40 and 10/40 passing predictions in the deposit.
  \item Our predictor is AlphaFold3 rather than Chai; best-of-$N$ uses 5
        diffusion samples per sequence.
  \item \textbf{Caveat to resolve before submission.} Our implementation of the
        one predictor-free benchmark quantity, \texttt{ligand\_dist\_des\_ncac\_min},
        does not reproduce the published values: median absolute error
        \SI{0.42}{\angstrom} over 40 backbones, \emph{zero} exact matches. Since
        it shares parsing and superposition code with the motif RMSD, the
        absolute values are not yet comparable to theirs. Until it is closed,
        only the backbone that reproduces (M0097) supports a head-to-head claim.
  \item \textbf{This is a metric-definition gap, not a reference-structure
        gap.} An earlier draft attributed it to our scoring against the
        idealized backbone where RFdiffusion2 scores against the unidealized,
        packed design. That explanation is now excluded by direct measurement.
        We recovered every candidate design structure from the authors'
        \SI{177}{\giga\byte} benchmark archive --- idealized, unidealized,
        LigandMPNN-packed, LigandMPNN input backbone --- and all four give the
        same distance to within \SI{0.01}{\angstrom} while all four miss the
        published value by ${\sim}\SI{0.42}{\angstrom}$. The authors' own
        \texttt{metrics\_cache} confirms the mechanism could not have worked:
        idealization displaces atoms by \SI{0.097}{\angstrom} mean and
        \SI{0.274}{\angstrom} max, an order of magnitude too small. Atom-subset
        and motif-inclusion variants are likewise excluded.
\end{planlist}
\todo{ame --- reproduce \texttt{ligand\_dist\_des\_ncac\_min} exactly; needs
  \texttt{per\_sequence\_metrics.py}, which is not in the deposit (outputs
  only). Structure hypothesis falsified 2026-08-15 --- do not retry it.}

\plan{\textbf{Scope of this caveat.} It bounds only whether our absolute
numbers may be set beside the published ones. Every comparison the paper
actually makes --- \cleo{} against LigandMPNN, the temperature ladder, the KL
sweep, the objective ablation --- runs through a single pipeline with a single
reference, and none of them is affected.}

\subsection*{Block 2 --- Low-temperature sampling returns one solution,
sampled repeatedly}

The incumbent protocol samples LigandMPNN at low temperature. At $n = 96$ on
M0097 this yields a high success rate --- \SI{58.3}{\percent} of sequences pass
--- but the passing designs are a tight cluster: mean pairwise Hamming distance
\num{40.6} out of $L = 180$, and \num{162} distinct substitutions among any 15 of
them. The library is close to one solution measured many times, and most of the
folding budget spent distinguishing its members buys little.

Raising the temperature removes the redundancy and most of the signal with it.
At $T = 1.0$ the same backbone reaches mean pairwise distance \num{126.3} at a
success rate of \SI{4.2}{\percent}. Diversity alone is not the objective; a
library that cannot be filtered is not a library.

\plan{Two numbers here were wrong in earlier drafts and are corrected. $T = 1.0$
was reported as \emph{zero} passing designs; that came from $n = 8$, and at
$n = 96$ it passes \SI{4.2}{\percent} of the time. And the low-temperature arm
was described as collapsing to a single cluster ``at every threshold down to
\SI{50}{\percent} identity'' --- true at \SI{70}{\percent}, but at
\SI{90}{\percent} every passing design is its own cluster and the statistic
simply counts designs. Cluster counts are threshold-dependent and are no longer
quoted; mean pairwise Hamming and the distinct-substitution count $U$ are
threshold-free and replace them throughout.}

\subsection*{Block 3 --- \cleo{} occupies the corner neither temperature reaches}

Sampling temperature traces a frontier: success falls as diversity rises, and
every point on it is available to the incumbent by turning one knob. The claim
of this section is that \cleo{} sits above that frontier, and the measurement is
in Table~\ref{tab:ame-frontier}.

\begin{table}[htbp]
\centering
\caption{M0097, per-sequence success rate on best-of-5 AF3 predictions, with
diversity measured over the \emph{passing} designs only --- a library contains
only what survives filtering. $U_{15}$ is the number of distinct
(position, residue) substitutions in a library of 15 passing designs, rarefied
over 300 draws so arms with different yields are comparable. Baselines are
$n = 96$; \cleo{} arms are the final 40 steps of a 200-step run.}
\label{tab:ame-frontier}
\begin{tabular}{@{}lcccc@{}}
\toprule
\textbf{Arm} & \textbf{Success} & \textbf{Mean pairwise} & $U_{15}$ &
\textbf{Median RMSD} \\
             &                  & \textbf{Hamming}       &          &
(\si{\angstrom}) \\
\midrule
LigandMPNN $T{=}0.1$ & \SI{58.3}{\percent} & 40.6  & 162 & 1.42 \\
LigandMPNN $T{=}0.2$ & \SI{66.7}{\percent} & 47.8  & 205 & 1.37 \\
LigandMPNN $T{=}0.3$ & \SI{58.3}{\percent} & 57.9  & 274 & 1.31 \\
LigandMPNN $T{=}0.5$ & \SI{34.4}{\percent} & 80.5  & 438 & 1.70 \\
LigandMPNN $T{=}0.7$ & \SI{26.0}{\percent} & 100.5 & 608 & 1.88 \\
LigandMPNN $T{=}1.0$ & \SI{4.2}{\percent}  & 126.3 & --- & 5.59 \\
\midrule
\cleo{}, no anchor         & \SI{83.1}{\percent} & 76.6 & 510 & 1.07 \\
\textbf{\cleo{}, KL anchor} & \textbf{\SI{56.7}{\percent}} & \textbf{114.4} &
  \textbf{808} & 1.44 \\
\bottomrule
\end{tabular}
\end{table}

Interpolating the baseline frontier at the anchored arm's diversity, the
incumbent succeeds \SI{14}{\percent} of the time where \cleo{} succeeds
\SI{56.7}{\percent} --- $\mathbf{3.98\times}$. Read the other way, at \cleo{}'s
success rate the incumbent reaches mean pairwise distance \num{59} against
\cleo{}'s \num{114.4} --- $\mathbf{1.93\times}$. The anchored arm also carries
more distinct substitutions per passing design than any baseline
($U_{15} = 808$ against \num{608} at $T{=}0.7$), so this is a genuine
exploration advantage and not only a higher yield.

The unanchored arm is the instructive contrast. It reaches the highest success
rate in the table, \SI{83.1}{\percent}, and the \emph{worst} per-design coverage
of any \cleo{} arm --- $U_{15} = 510$, below $T{=}0.7$. Left unconstrained the
policy converges onto a narrow solution. The KL weight against the base model is
what holds the distribution open, and it is the control the library application
needs.

Sweeping that weight over $\{0, 0.005, 0.01, 0.02, 0.05, 0.1\}$ shows it is a
well-behaved control with an interior optimum rather than a knob that has to be
tuned. Success falls monotonically, \SIrange{83}{10}{\percent}, and per-design
coverage rises monotonically, $U_{15}$ \numrange{510}{881}; the total distinct
substitutions a library actually contains is the product of the two and is flat
within noise across $w \in [0.01, 0.05]$ (\numrange{2480}{2501}) before
collapsing to \num{1743} at $w = 0.1$, where too few designs pass for
per-design exploration to accumulate. The value used throughout,
$w = \num{0.02}$, sits in the middle of that plateau, so the reported numbers do
not depend on having found a sharp optimum.

% Shared figure block, \input by BOTH main.tex (via results/frontier.tex) and
% notebook.tex (via experiments.tex). The two are independent builds, so one
% label here cannot clash; keeping a single copy means the caption cannot drift
% between the paper and the record, which is exactly what happened to the CLEO
% frontier row before the split.
\begin{figure}[htbp]
  \centering
  \includegraphics[width=\textwidth]{ame_distributions}
  \caption{\textbf{Where the distributions actually sit, rather than the success
  rate thresholded out of them.} Each row is one arm's per-design best-of-5 motif
  RMSD: box for median and quartiles, jittered strip behind it for the shape the
  box hides. Log axis, because a linear one lets the single worst arm compress
  everything that matters into the left fifth. \textbf{Left, M0097:} every
  \cleo{} arm has its \emph{median} below the cutoff (\SIrange{1.07}{1.44}{\angstrom})
  --- the mode has moved, not just the tail --- against
  \SIrange{1.31}{1.42}{\angstrom} for low-temperature LigandMPNN and
  \SI{5.59}{\angstrom} at $T{=}1.0$. The non-monotonicity noted in E10 is visible
  directly: $T{=}0.3$ has a \emph{lower} median than $T{=}0.1$.
  \textbf{Right, the three near-miss backbones:} medians of
  \SIrange{2.81}{3.12}{\angstrom} with the cutoff far out in the left tail. The
  one design that survived 40 refolds sits at the extreme edge of its own
  distribution. Encoding: both ladders are ordinal, so each is a single hue
  ramped light to dark rather than arbitrary categorical colours; the anchored
  arm keeps an accent hue as the highlighted result.}
  \label{fig:ame-distributions}
\end{figure}


Figure~\ref{fig:ame-distributions} shows why the success rate alone is not
enough: on M0097 every \cleo{} arm has its \emph{median} below the
\SI{1.5}{\angstrom} cutoff, so the mode has moved rather than the tail merely
reaching across --- which is precisely what fails to hold on the near-miss
backbones of Block 6.

\subsection*{Block 4 --- The training trajectory is the library}

\plan{This is the reframing that removes a hyperparameter. Rather than adding an
explicit diversity term to the reward, we treat every rollout sampled during
training as a library candidate. Those sequences were already folded and scored
to compute the reward, so the candidate pool is free: the library is a byproduct
of training, not an additional expense.}

Across 200 GRPO steps at 16 sequences per step, M0097 accumulates 3200 scored
designs. Under the anchored policy \num{1349} of them are distinct passing
sequences carrying \num{3045} distinct substitutions --- \SI{89}{\percent} of the
$19L = 3420$ ceiling.

\begin{center}
\begin{tabular}{@{}lcccc@{}}
\toprule
& \multicolumn{2}{c}{\textbf{\cleo{}, no anchor}} &
  \multicolumn{2}{c}{\textbf{\cleo{}, KL anchor}} \\
\cmidrule(lr){2-3}\cmidrule(lr){4-5}
\textbf{Cumulative folds} & \textbf{Unique passing} & \textbf{Substitutions} &
\textbf{Unique passing} & \textbf{Substitutions} \\
\midrule
400  & 75   & 1483 & 47   & 1382 \\
800  & 284  & 2294 & 160  & 2045 \\
1600 & 893  & 2805 & 510  & 2757 \\
2400 & 1560 & 2974 & 903  & 2958 \\
3200 & 2229 & 3040 & 1349 & 3045 \\
\bottomrule
\end{tabular}
\end{center}

The distinct-substitution count is genuinely saturating here, at
\SIrange{89}{89}{\percent} of the ceiling, while the count of distinct passing
sequences is still climbing steeply. The library is not running out of designs;
it is running out of \emph{single} substitutions to explore, which is the
expected behaviour once most positions have been visited and further diversity
must come from combinations.

\plan{This supersedes an earlier version of this table reporting 97 passing
designs from 3200 folds, a factor of \num{14} to \num{23} lower. That run used
the pre-audit objective, and its curve flattened because the policy decayed
rather than because sequence space was exhausted --- the earlier draft said so
and predicted rank normalisation would keep it climbing, which is what happened.
The pool numbers here are best-of-5, matching Block 3; the superseded ones were
single-sample and were not comparable to it.}

\subsection*{Block 4b --- Where in sequence space the passing designs sit}

\plan{Generated by \texttt{paper/figures/ame\_pca.py}, which writes two figures
because one projection cannot honestly carry both claims. The shared projection
is fit across arms within a backbone --- a per-arm PCA places every cloud at the
origin and destroys the comparison --- but deliberately not across backbones,
where a global fit puts \SI{99}{\percent} of PC1 variance on backbone identity.

\textbf{These two panels are stale and must be regenerated.} They were built
from the pre-audit runs and from baseline arms holding 8 sequences each. Both
inputs have been replaced: the objective is fixed and the baselines are now
$n = 96$. The drift control below remains the right control and its conclusion
is unlikely to change, but no number in these captions should be quoted until
they are rebuilt.

Explained variance is low (PC1 \SIrange{5}{9}{\percent}), as expected for
one-hot protein sequence, so the projection is a visual aid; every quantitative
diversity claim in this section uses Hamming distance and substitution counts
directly.}

\begin{figure}[htbp]
  \centering
  \includegraphics[width=\textwidth]{ame_pca}
  \caption{\textbf{Sequence space occupied by each design method on the AME
  pilot backbones.} One PCA per backbone, fit on every sequence that backbone
  contributes across all three arms, so the arms share one projection and the
  clouds are directly comparable. Faint points are sampled designs, solid
  points are designs that pass the benchmark criterion (motif all-atom RMSD
  $<$ \SI{1.5}{\angstrom}, no ligand clash). \cleo{} is restricted to its
  peak-performing 10-step window so that three \emph{policies} are compared at
  one point in training. On M0097 the low-temperature arm (orange) forms a
  tight knot displaced far along PC1 --- five passing designs that are a single
  cluster at every identity threshold down to \SI{50}{\percent} --- while
  \cleo{}'s passing designs are spread throughout a broad cloud. $T{=}1.0$
  (grey) reaches comparable spread with zero passing designs.
  Projections are deliberately \emph{not} shared across backbones: a global fit
  places \SI{99}{\percent} of PC1 variance on backbone identity, collapsing
  each backbone to its own blob and hiding the arm differences entirely.
  Explained variance is low (PC1 \SIrange{5}{9}{\percent}), as expected for
  one-hot encoded protein sequence; the projection is a visual aid and every
  quantitative claim in the text uses Hamming distance and cluster counts
  directly. Baseline arms hold 8 sequences each, so part of their compactness
  is small-$n$.}
  \label{fig:ame-pca}
\end{figure}

\begin{figure}[htbp]
  \centering
  \includegraphics[width=\textwidth]{ame_pca_drift}
  \caption{\textbf{Most of PC1 is drift over training, not diversity available
  at any one step.} The full GRPO trajectory for each backbone, in the same
  projection as Figure~\ref{fig:ame-pca}, coloured by training step; passing
  designs are circled. PC1 correlates with training step at $r = +0.92$ to
  $+0.94$ on all three backbones. This distinction is load-bearing rather than
  cosmetic: spread accumulated by drifting across training is genuine library
  diversity when the library is assembled from the whole run, which is how ours
  is assembled, but it is \emph{not} evidence that any single policy is
  diverse. Presenting the pooled projection alone would conflate the two and
  overstate the result, so the axis is shown rather than hidden.}
  \label{fig:ame-pca-drift}
\end{figure}

\todo{ame --- rebuild both PCA panels on the corrected runs at matched $n$}

\subsection*{Block 5 --- The temperature frontier}

Measured, and reported in Table~\ref{tab:ame-frontier}. Two predictions stated
in advance were tested.

\textbf{The dangerous intermediate temperatures do not close the gap.} $T=0.2$
and $T=0.3$ were named as the points that would refute a frontier shift if
either reached \cleo{}'s coverage at \cleo{}'s success rate. Neither does: both
sit below \num{0.6} of the anchored arm's diversity. The frontier shift stands.

\textbf{One prediction was wrong.} The trade-off was expected to be monotone in
temperature; it is not. $T=0.2$ is the \emph{most} successful arm in the table
at \SI{66.7}{\percent}, above $T=0.1$'s \SI{58.3}{\percent}, while also being
more diverse. $T=0.2$, not $T=0.1$, is therefore the incumbent to quote at the
low-diversity end, and any comparison against $T=0.1$ alone flatters us.

\plan{Still outstanding from this block's original plan: the saturation fold
count for low-temperature LigandMPNN --- the point past which it re-emits
sequences it has effectively already produced. The \cleo{} side of that
comparison is Block 4; the baseline side needs a cumulative-distinct curve at
matched budget, which the $n = 96$ sweep can supply but does not yet.}

\subsection*{Block 6 --- Rescue of backbones the benchmark cannot solve}

Three backbones sit at \SIrange{1.61}{1.69}{\angstrom} best-of-40 in the deposit
--- zero passing predictions, but only ${\sim}\SI{0.15}{\angstrom} $ from
threshold. A pass on any of them is a design the published pipeline cannot
produce, with no reproduction confound, because there is no positive result of
theirs for us to fail to reproduce.

All three were trained for 75 steps at batch 16. Nine designs cleared the
criterion. \textbf{Each was then refolded 40 times --- the published
per-backbone budget --- and one survives.} M0365 \texttt{hit00} passes
\num{22} of \num{40} predictions with a \emph{median} RMSD of
\SI{1.470}{\angstrom}, below the cutoff rather than touching it, on a backbone
the deposit scores 0/40 with a best of \SI{1.69}{\angstrom}. That is a design
the published pipeline does not produce.

The other eight do not survive. Five pass zero of 40; three pass one or two.
Training-time margin predicts this almost perfectly: the single design with a
margin above \SI{0.03}{\angstrom} inside the cutoff is the single design that
replicates, and everything at \SIrange{0.003}{0.023}{\angstrom} was a tail draw
from a distribution whose median sits near \SI{3}{\angstrom}
(Figure~\ref{fig:ame-distributions}, right).

\plan{\textbf{The honest claim is one rescue, not nine}, and the general lesson
is worth more than the rescue: a best-of-$N$ filter applied to a distribution
whose mode is far from the cutoff selects excursions, not solutions. A design
that truly passes 1 in 40 times is flagged by a best-of-5 filter about
\SI{12}{\percent} of the time it is folded. Every headline number that leaves
this section should be replicated the way \texttt{hit00} was, and in-training
pass rates should be described as counts of best-of-$N$ events rather than of
reproducible designs.

Two caveats remain open. The fold budget is not matched --- 1200 designs at
best-of-5 against the deposit's 8 sequences at 5 models --- so a compute-matched
comparison is still owed. And the idealised-versus-unidealised reference
structure question from Block 1 applies here as it does everywhere.}
